\documentclass[12pt]{article}
\usepackage{amsmath,amssymb}

\begin{document}
\section*{{2023 AIME I Problems/Problem 11}}
Source: AoPS Wiki

\subsection*{Solution}
Define $f(x)$ to be the number of subsets of $\{1, 2, 3, 4, \ldots x\}$ that have $0$ consecutive element pairs, and $f'(x)$ to be the number of subsets that have $1$ consecutive pair. Using casework on where the consecutive element pair is, there is a unique consecutive element pair that satisfies the conditions. It is easy to see that \[f'(10) = 2f(7) + 2f(6) + 2f(1)f(5) + 2f(2)f(4) + f(3)^2.\] We see that $f(1) = 2$ , $f(2) = 3$ , and $f(n) = f(n-1) + f(n-2)$ . This is because if the element $n$ is included in our subset, then there are $f(n-2)$ possibilities for the rest of the elements (because $n-1$ cannot be used), and otherwise there are $f(n-1)$ possibilities. Thus, by induction, $f(n)$ is the $n+1$ th Fibonacci number. This means that $f'(10) = 2(34) + 2(21) + 2(2)(13) + 2(3)(8) + 5^2 = \boxed{235}$ . ~mathboy100

\end{document}
